\chapter{Fairness, Explainability, and Governance}
\label{chap:fairness}

\section{SHAP Explainability}
\label{sec:shap}

SHAP (SHapley Additive exPlanations)~\cite{lundberg2017} TreeExplainer was applied to all four model families across both traditional and combined datasets, producing portfolio-level feature importance rankings and individual-level prediction explanations. SHAP values have the desirable property of being theoretically grounded in cooperative game theory: each feature's contribution to a prediction is computed as its average marginal contribution across all possible feature coalitions, ensuring consistency and local accuracy.

\subsection{Traditional vs Combined SHAP Comparison}

The SHAP analysis was conducted for each model family on both the traditional (application-only) and combined (application + supplementary) datasets to understand how the addition of alternative data changes the model's decision-making logic.

\subsubsection{Logistic Regression}

\ProjectTwoFigure[!htbp][width=0.45\textwidth][width=0.45\textwidth]{pptx/SHAP explainability - logistic regression - traditional.png}{pptx/SHAP explainability - logistic regression - combined.png}{SHAP LR: Trad vs Combined}{SHAP feature importance for Logistic Regression: Traditional dataset (left) vs Combined dataset (right). The traditional model relies heavily on employment duration and income type, while the combined model shifts toward external source scores.}{fig:shap_lr}


In the traditional dataset, a few features---especially \texttt{DAYS\_EMPLOYED}, income type, \texttt{AMT\_CREDIT}, and the \texttt{EXT\_SOURCE} variables---explain most of the variation in predicted default risk. The concentration of importance in a small number of features is characteristic of linear models, which allocate predictive weight through a single coefficient per feature. In the combined dataset, Logistic Regression shifts toward external sources as the dominant drivers, with behavioural features providing marginal additional signal.

\ProjectFigure[!htbp][width=0.62\textwidth]{notebook/final_traditional_shap_summary_plot_bar_shows_fe_3.png}{SHAP LR Detail}{Detailed SHAP importance bar plot for Logistic Regression on the traditional dataset. The dominant features align with the paired traditional-versus-combined comparison above, so the repeated beeswarm panel is omitted here for concision.}{fig:shap_lr_detailed}


\subsubsection{Random Forest}

\ProjectTwoFigure[!htbp][width=0.45\textwidth][width=0.45\textwidth]{pptx/SHAP explainability - random forest - traditional.png}{pptx/SHAP explainability - random forest - combined.png}{SHAP RF: Trad vs Combined}{SHAP feature importance for Random Forest: Traditional (left) vs Combined (right). External scores dominate in both datasets, with other features contributing substantially less.}{fig:shap_rf}


Random Forest models are primarily driven by external scores in both datasets. Age (\texttt{DAYS\_BIRTH}), education type, and financial ratios contribute, but their overall impact is much smaller relative to the \texttt{EXT\_SOURCE} variables. The Random Forest's bagging mechanism distributes importance more evenly across correlated features compared to boosting methods, resulting in a flatter importance tail.

\ProjectFigure[!htbp][width=0.62\textwidth]{notebook/final_traditional_shap_summary_plot_bar_shows_fe_6.png}{SHAP RF Detail}{Detailed SHAP importance bar plot for Random Forest on the traditional dataset. EXT\_SOURCE variables remain dominant, with the importance tail flatter than the boosting models due to bagging.}{fig:shap_rf_detailed}


\subsubsection{XGBoost}

\ProjectTwoFigure[!htbp][width=0.45\textwidth][width=0.45\textwidth]{pptx/SHAP explainability - xgboost - traditional.png}{pptx/SHAP explainability - xgboost - combined.png}{SHAP XGB: Trad vs Combined}{SHAP feature importance for XGBoost: Traditional (left) vs Combined (right). The combined model incorporates a wider variety of engineered features, spreading predictive influence across more diverse financial behaviours.}{fig:shap_xgb}


The traditional XGBoost model relies heavily on external source variables and \texttt{AMT\_GOODS\_PRICE}. The combined model shows a notable shift: it incorporates a wider variety of engineered features such as \texttt{CREDIT\_ANNUITY\_PERCENT} and instalment history, spreading predictive influence across more diverse financial behaviours. This diversification is beneficial for model robustness---a model that relies less on any single feature is less vulnerable to feature drift or data quality issues.

\ProjectFigure[!htbp][width=0.62\textwidth]{notebook/final_traditional_shap_summary_plot_bar_shows_fe_10.png}{SHAP XGB Detail}{Detailed SHAP importance bar plot for XGBoost on the traditional dataset. The sharper drop-off after the top-ranked features reflects the more concentrated gradient-boosting structure.}{fig:shap_xgb_detailed}


\subsubsection{LightGBM}

\ProjectTwoFigure[!htbp][width=0.45\textwidth][width=0.45\textwidth]{pptx/SHAP explainability - LightGBM - traditional.png}{pptx/SHAP explainability - LightGBM - combined.png}{SHAP LGBM: Trad vs Combined}{SHAP feature importance for LightGBM: Traditional (left) vs Combined (right). The combined model leverages aggregated data from previous applications and instalment payment history, producing the most diversified feature importance profile.}{fig:shap_lgbm}


LightGBM shows similar patterns to XGBoost due to their shared gradient-boosting foundations. The combined model leverages more aggregated data from previous applications and instalment payment history. LightGBM's leaf-wise growth strategy tends to create deeper, more specialised trees that can extract value from a broader set of features, resulting in the most diversified feature importance profile among all four models.

\ProjectFigure[!htbp][width=0.62\textwidth]{notebook/final_traditional_shap_summary_plot_bar_shows_fe_14.png}{SHAP LGBM Detail}{Detailed SHAP importance bar plot for LightGBM on the traditional dataset. The diversified importance profile is consistent with the champion model's broader use of behavioural and contextual features.}{fig:shap_lgbm_detailed}


\subsection{Cross-Model Feature Importance Summary}

Across all four models and both datasets, three consistent patterns emerge:

\begin{enumerate}
    \item \textbf{EXT\_SOURCE dominance:} \texttt{EXT\_SOURCE\_2}, \texttt{EXT\_SOURCE\_3}, and \texttt{EXT\_SOURCE\_1} consistently rank as the top three features, regardless of model family or dataset variant. Their combined SHAP importance accounts for approximately 40\% of total model signal.
    \item \textbf{Age as a universal predictor:} \texttt{DAYS\_BIRTH} (age) appears in the top 5--10 features for every model, with a consistent directional effect: younger applicants have higher predicted default risk.
    \item \textbf{Combined models diversify:} When supplementary data is added, the importance of application-level features decreases as behavioural features absorb some of the predictive signal. This is desirable from a robustness perspective, as it reduces single-feature dependency.
\end{enumerate}

The paired SHAP comparisons already show the dominant risk drivers clearly, so the previously duplicated miscaptioned figure is removed from this chapter and reused in Chapter~\ref{subsec:prev_app_eda}, where it actually belongs as a previous-application contract-status visual.

\section{LIME Explanations}
\label{sec:lime}

While SHAP provides theoretically exact attributions for tree-based models, LIME (Local Interpretable Model-agnostic Explanations)~\cite{ribeiro2016} was applied to provide complementary individual-level prediction explanations. LIME operates by fitting a local linear model around each prediction, offering an intuitive ``which features pushed this decision and by how much'' explanation.

For each major model family, one default and one non-default example were explained, with local feature contributions visualised. The LIME explanations consistently agreed with the SHAP global rankings for the top features, reinforcing confidence in the model's decision logic. Where LIME and SHAP disagreed on feature ordering for lower-ranked features, this reflected LIME's local approximation versus SHAP's global exact computation---for regulatory purposes, SHAP is preferred for its theoretical guarantees, while LIME is useful for generating intuitive customer-facing explanations.

\section{Fairness Evaluation}
\label{sec:fairness_eval}

Fairness evaluation was conducted using the Fairlearn library~\cite{bird2020}, applying multiple fairness definitions to assess whether the model's predictions are equitable across demographic groups.

\subsection{Gender Fairness Analysis}

Two primary metrics were used to quantify potential gender-based bias:

\begin{itemize}
    \item \textbf{Demographic Parity (Selection Rate):} Measures whether each gender group has an equal probability of being predicted as default. The demographic parity difference is $|P(\hat{Y}=1|G=\text{male}) - P(\hat{Y}=1|G=\text{female})|$. Values closer to 0 indicate less bias.
    \item \textbf{Equalized Odds (Recall Disparity):} Checks whether the model is equally accurate at catching actual defaults across groups~\cite{hardt2016}. The equalized odds difference is $|TPR_\text{male} - TPR_\text{female}|$. A model satisfying equalized odds makes errors at equal rates across groups.
\end{itemize}

\ProjectFigure[!htbp][width=0.8\textwidth]{pptx/fairness evaluation by gender.png}{Fairness by Gender}{Fairness evaluation by gender across all four model families. Demographic parity disparity varies by model, while equalized odds disparity is uniform---indicating that bias originates from the data, not the algorithm.}{fig:fairness_gender}


\ProjectFigure[!htbp][width=0.85\textwidth]{notebook/final_traditional_visualize_fairness_metrics_acr_16.png}{Traditional Fairness Metrics}{Detailed fairness metrics for the traditional dataset across all models, showing selection rates, true positive rates, and false positive rates by gender group.}{fig:fairness_traditional}


Key findings:

\begin{enumerate}
    \item \textbf{Random Forest shows the lowest demographic parity disparity}, making the most consistent predictions across genders. Its bagging mechanism distributes decision-making across many trees, diluting the effect of any single gender-correlated feature.
    \item \textbf{Logistic Regression has the highest disparity}, suggesting the strongest imbalance in outcomes between genders. Linear models absorb data biases directly into their coefficients, amplifying any correlation between gender-related features and the target.
    \item \textbf{All four models have uniform equalised odds}, meaning the bias is in the dataset itself, not in the algorithm of choice. The correct intervention is therefore upstream---resampling, reweighting, or collecting more representative training data.
    \item \textbf{Disparate impact ratios remain $\geq 0.80$} across all models, meeting the four-fifths rule threshold commonly used in regulatory compliance assessments.
\end{enumerate}

\subsection{Gender Bias Analysis: The XNA Group}

\ProjectFigure[!htbp][width=0.8\textwidth]{pptx/gender bias analysis.png}{Gender Bias by Group}{Selection rates by gender across all models. The XNA group (non-disclosed gender) reveals a critical failure mode for Logistic Regression, which predicts exactly 0 for this group.}{fig:gender_bias}


The XNA group---applicants who chose not to disclose their gender (4 records in the dataset)---reveals an important failure mode. Logistic Regression predicts exactly 0 default probability for this group, a statistically impossible result caused by the model's inability to learn from such a small sample size. This is a dangerous overconfident prediction: it would auto-approve all non-disclosing applicants regardless of their actual risk profile.

Ensemble models (Random Forest, XGBoost, LightGBM) revert to the base rate ($\sim$0.50) for the XNA group, expressing appropriate uncertainty rather than making a dangerously confident prediction. This behaviour arises because ensemble models average across many trees, each of which may handle the rare group differently, producing a prediction closer to the population prior.

\subsection{The Accuracy--Fairness Trade-off}

Logistic Regression achieves the highest overall recall but shows the worst demographic parity and fails catastrophically on the XNA group. Its single global linear boundary aggressively exploits all correlations, including gender-correlated signals such as occupation type and income patterns.

Ensemble models sacrifice a small amount of aggregate accuracy in exchange for more distributed and balanced decision-making across all customer segments. Since this project is designed to serve thin-file customers who are already marginalised by traditional credit systems, deploying the highest-recall model without fairness evaluation would be exactly the wrong decision.

This trade-off motivates the argument for preferring ensemble models (specifically LightGBM) for final predictions: they are more robust when handling underrepresented customer segments, even if a linear model achieves slightly better recall on the majority population.

\subsection{Bias Mitigation Strategies}

Three complementary mitigation strategies were explored:

\begin{enumerate}
    \item \textbf{ThresholdOptimizer (post-processing):} Fairlearn's \texttt{ThresholdOptimizer}~\cite{bird2020} was applied under demographic parity constraints to equalise selection rates across gender groups. This approach adjusts decision thresholds per group without retraining the model, preserving the model's learned patterns while correcting for disparate outcomes. The AUC-ROC impact was minimal ($< 0.005$).

    \item \textbf{Feature masking:} Sensitive attributes (gender) were monitored but not directly used as prediction features. The model was trained to find patterns in credit history and financial ratios independent of gender. However, proxy features (e.g., occupation type, income level) can still carry gender-correlated signal, making complete debiasing through feature removal alone insufficient.

    \item \textbf{Group-specific threshold optimisation:} Rather than using a uniform 0.5 cutoff for all groups, group-specific thresholds were evaluated to equalise selection rates. This achieves demographic parity by construction but may sacrifice equalized odds---a trade-off that must be evaluated against the specific regulatory framework applicable to the lending jurisdiction.
\end{enumerate}

The originally inserted ``before/after fairness'' graphics were in fact IV feature-selection artifacts, so they have been relocated to Chapter~\ref{sec:feature_selection}. This chapter therefore keeps the fairness discussion on genuinely fairness-specific visuals and measured disparities.

\section{Model Monitoring Framework}
\label{sec:monitoring}

A comprehensive model monitoring framework was developed following SR 11-7 guidance on model risk management~\cite{occ2011}. The framework tracks four dimensions of model health: population stability, feature stability, discriminatory power, and fairness stability.

\subsection{Population Stability Index (PSI)}

PSI~\cite{yurdakul2020} measures the shift in the score distribution between the development (training) and monitoring (validation/production) populations:

\begin{equation}
PSI = \sum_{i=1}^{n} (p_i^{\text{actual}} - p_i^{\text{expected}}) \times \ln\left(\frac{p_i^{\text{actual}}}{p_i^{\text{expected}}}\right)
\end{equation}

The industry-standard thresholds are:
\begin{itemize}
    \item PSI $< 0.10$: No significant drift; model is stable
    \item $0.10 \leq$ PSI $< 0.25$: Moderate drift; investigation warranted
    \item PSI $\geq 0.25$: Significant population shift; model retraining required
\end{itemize}

\ProjectFigure[!htbp][width=0.85\textwidth]{notebook/home_credit_modeling_combined_psi_distribution_visualization_28.png}{PSI Distribution}{PSI distribution visualisation comparing score distributions between training and validation populations. The PSI value confirms acceptable population stability.}{fig:psi}


\subsection{Characteristic Stability Index (CSI)}

CSI applies the same divergence formula as PSI but at the individual feature level, enabling identification of which specific features are driving any observed population shift. Features with CSI $> 0.25$ are flagged for investigation, as they may indicate changes in data collection, applicant demographics, or market conditions that could degrade model performance.

\subsection{AUC and Gini Stability}

AUC and Gini coefficient stability were monitored across time-based validation splits created using quarterly windows via the \texttt{DAYS\_REGISTRATION} variable. This temporal partitioning simulates how the model would perform as the applicant population evolves over time. Stable AUC across quarters (variation $< 0.02$) indicates that the model's discriminatory power is not dependent on temporal artefacts in the training data.

\subsection{Fairness Stability}

Fairness metrics---specifically the change in disparate impact ratio ($\Delta$DI $< 0.02$)---were tracked across validation splits to ensure that fairness properties are maintained over time, not just at a single evaluation point. A model that is fair on the development sample but drifts toward unfairness on new populations provides a false sense of compliance.

\section{Model Governance}
\label{sec:governance}

\subsection{Regulatory Compliance Framework}

The model governance framework was designed with Home Credit's multi-country operations in mind, addressing regulatory requirements across Southeast Asian jurisdictions:

\begin{itemize}
    \item \textbf{Consent Framework:} Opt-in data sharing with granular permissions per feature category. Applicants can consent to sharing application data while declining behavioural data sharing, with the model gracefully degrading to the traditional-only variant.
    \item \textbf{Anonymisation:} $k$-anonymity ($k \geq 5$) for demographic data and differential privacy ($\epsilon = 1.0$) for spending and financial data, ensuring that individual records cannot be re-identified from aggregated model outputs.
    \item \textbf{Data Encryption:} AES-256 encryption at rest and TLS 1.3 in transit for all applicant data, protecting against unauthorised access and cyber attacks.
    \item \textbf{Data Residency:} In-country data storage compliant with PDPA (Thailand), PDPA (Singapore), and DPA (Philippines), ensuring that applicant data does not cross jurisdictional boundaries without explicit consent.
    \item \textbf{Right to Erasure:} Applicants can revoke data usage at any time. Upon erasure request, their data is removed from training sets and the model is scheduled for quarterly retraining to eliminate any residual influence.
\end{itemize}

\subsection{Risk Mitigation Matrix}

\begin{table}[H]
\centering
\caption{Risk Mitigation Matrix for Model Deployment}
\label{tab:risk_mitigation}
\begin{tabular}{p{3.5cm}p{2cm}p{2cm}p{5.5cm}}
\toprule
\textbf{Risk} & \textbf{Likelihood} & \textbf{Impact} & \textbf{Mitigation Strategy} \\
\midrule
Regulatory non-compliance & Medium & High & Periodic compliance checks; SHAP-based audit trail for every decision; jurisdiction-specific threshold configurations \\
Model bias against underrepresented groups & Medium & High & Quarterly fairness audits using Fairlearn MetricFrame; dynamic undersampling/oversampling; ThresholdOptimizer post-processing \\
Data partner dependency / API outage & Low & Medium & Multiple redundant data sources; graceful fallback to application-only model when external scores unavailable \\
Stakeholder hesitation / slow adoption & Medium & Medium & Transparent explainability dashboard; human advisor override capability; phased rollout with A/B testing \\
Population drift degrading model accuracy & Medium & High & Monthly PSI/CSI monitoring; automated retraining triggers at PSI $> 0.15$; champion-challenger framework \\
\bottomrule
\end{tabular}
\end{table}

\subsection{Model Documentation and Audit Trail}

Following SR 11-7 requirements~\cite{occ2011}, the model documentation package includes:

\begin{enumerate}
    \item \textbf{Model development documentation:} Feature engineering rationale, model selection criteria, hyperparameter choices, and validation results.
    \item \textbf{Ongoing monitoring plan:} PSI/CSI thresholds, AUC stability criteria, fairness metric bounds, and escalation procedures.
    \item \textbf{Decision audit trail:} SHAP values for every prediction enable post-hoc explanation of individual lending decisions, satisfying right-to-explanation requirements.
    \item \textbf{Change management protocol:} Model retraining, feature changes, and threshold adjustments require documented approval and validation before deployment.
\end{enumerate}
